\documentclass[10pt,a4paper]{article}
\usepackage[T1]{fontenc}
\usepackage[utf8]{inputenc}
\usepackage[english]{babel}
\usepackage{lmodern}
\usepackage[margin=1.65cm]{geometry}
\usepackage{booktabs,tabularx,array}
\usepackage[hidelinks]{hyperref}
\setlength{\parindent}{0pt}
\setlength{\parskip}{4pt}
\renewcommand{\arraystretch}{1.12}
\newcommand{\heading}[1]{\vspace{5pt}{\large\bfseries #1}\par\vspace{2pt}}
\newcommand{\code}[1]{\texttt{\small #1}}
\pagestyle{plain}

\begin{document}
\begin{center}
{\LARGE\bfseries Classifying reported cardiovascular disease\par}
\vspace{3pt}
{\large BRFSS 2015: choices, experiments and results}\par
\vspace{3pt}
Antonio Arbizzani Tonelli, Alessia Meroni, Luca Sabatini\par
\vspace{3pt}
EPFL CS-433 --- Project 1
\end{center}
\vspace{-3pt}

We selected Phase 14 histogram gradient boosting implemented in NumPy.
Its AIcrowd submission obtains \textbf{F1 0.452 and accuracy 0.878}.
Nested development F1 is 0.44216, compared with 0.42563 for logistic regression.
Codebook-based cleaning and subsequent experiments led us to keep 295 features,
direct handling of missing values, all training rows and 200 trees.

\heading{Data and evaluation}
The dataset contains 328,135 labelled rows, 321 predictors and 109,379 test
rows. The target \code{\_MICHD} records self-reported myocardial infarction
or coronary heart disease. Positives make up 8.83\% of training labels;
44.79\% of predictor cells are blank. We used F1 as our main criterion:
always predicting negative gives 91.17\% accuracy and F1 zero.

The initial stratified split uses seed 20260918: 262,508 development rows and
65,627 local validation rows. The reported comparisons use three outer folds
on development data (seed 20260919). Within each outer training partition,
two inner folds produce out-of-fold predictions and select the threshold
maximising F1. We then apply that threshold to the excluded outer fold.
Preprocessing, imputation, feature selection and sampling use only the training
partition of each fit. Validation retains all rows and its original prevalence.
Pooled F1 uses the outer predictions, with the threshold
chosen separately for each fold.

\heading{Methods and preprocessing choices}
In \code{implementations.py}, we implemented all six required functions:
linear regression with GD/SGD, least squares, ridge, logistic regression and
L2 logistic regression. Least squares and ridge solve the normal equations;
linear losses use half-MSE. Returned losses exclude regularisation.

Linear baselines use historical median imputation, training-fitted
standardisation and age interactions. Logistic regression starts from zero,
with $\gamma=0.2$, 4,000 updates and $\lambda=0$.
Ridge uses $\lambda=10^{-4}$, selected from three values, with logistic
calibration fitted on inner OOF scores. These results precede categorical
mode imputation.

\textbf{Codebook-based preprocessing.} Using the
\href{https://www.cdc.gov/brfss/annual_data/2015/pdf/CODEBOOK15_LLCP.pdf}{CDC 2015 codebook},
we audited types, units, response codes and redundancies. We removed identifiers, the full interview
date, survey-design fields and redundant measurements, keeping 295 source
columns. Nonresponse codes become \code{NaN} only where the codebook assigns
that meaning; \code{ALCDAY5=888} and \code{PHYSHLTH=88} become zero.
This preserves valid category meanings and separates nonresponse from zero
quantities. Across the retained development columns, we recode 1.10 million
finite nonresponses and about 900,000 documented zero quantities.
Questionnaire branching motivates retaining a distinct missing
state: each tree split learns its direction from training rows.
Quantile bins are also fitted on training data. Each new tree uses the
gradient and curvature of logistic loss to update predictions.

An implementation audit found omitted splits separating all observed
values from missing responses. Phase 34 restores this candidate and was
evaluated on the same features, parameters and nested folds. The submitted
model preserves the historical Phase 14 search for exact reproduction.

The selected configuration uses \textbf{200 trees, depth 5}, learning rate
0.05, at least 200 rows per leaf, 64 bins, 128 candidate features per node,
leaf L2 of 1, positive weight 1 and seed 20260920.
The initial grid compared depths 3/5, 100/200/400 trees and 64/128 candidates.
Before the split correction, additional codebook corrections raise F1 from
0.44174 to 0.44216, with mixed directions across folds. Their consistent
interpretation of response codes supports retaining them.

\heading{Results on the same development data}
\begin{center}
\small
\begin{tabularx}{\linewidth}{@{}Xrrrr@{}}
\toprule
Pipeline & F1 & Precision & Recall & AP \\
\midrule
Ridge & 0.41951 & 0.33303 & 0.56665 & 0.39028 \\
Logistic regression & 0.42563 & 0.35728 & 0.52632 & 0.38775 \\
Raw booster, earlier split search & 0.44146 & 0.36663 & 0.55466 & 0.42783 \\
\textbf{Selected Phase 14 booster} & \textbf{0.44216} & \textbf{0.36763} & \textbf{0.55457} & \textbf{0.42861} \\
Phase 34, complete split search & 0.44080 & 0.36200 & 0.56346 & 0.42878 \\
MLP 64/32 & 0.42445 & 0.36072 & 0.51553 & 0.38778 \\
\bottomrule
\end{tabularx}
\end{center}
\vspace{-4pt}
The selected booster improves F1 by \textbf{0.01653} over logistic regression, increasing
both precision and recall. Average precision (AP) also supports better
ranking of positives. The three fold F1 scores are 0.43903, 0.44405 and 0.44351.

\newpage
\heading{Experiments that guided the final choice}
The control is historical Phase 14 (F1 0.44216), except $\dagger$ rows, which
use corrected Phase 34 (0.44080). All thresholds are selected on inner folds;
$\Delta$ is the difference from the corresponding control.
\begin{center}
\small
\begin{tabularx}{\linewidth}{@{}Xrr@{}}
\toprule
Experiment & F1 & $\Delta$F1 \\
\midrule
Daily food frequencies / activity one-hot & 0.44186 / 0.44057 & $-0.00029$ / $-0.00159$ \\
Exact discrete bins / eight nonresponse indicators & 0.44078 / 0.44032 & $-0.00137$ / $-0.00183$ \\
Selective imputation & 0.44179 & $-0.00036$ \\
100 / 60 features selected by permutation & 0.43922 / 0.43758 & $-0.00293$ / $-0.00457$ \\
Learning rate 0.025, 400 trees / depth 7 & 0.44059 / 0.44116 & $-0.00157$ / $-0.00099$ \\
Positive weight 1.25 / 1.5 & 0.44083 / 0.44061 & $-0.00132$ / $-0.00154$ \\
Mean probabilities across two random seeds & 0.44111 & $-0.00104$ \\
Removal of \code{HAREHAB1} & 0.43104 & $-0.01112$ \\
Random undersampling 5:1 / 7:1 & 0.43823 / 0.43982 & $-0.00393$ / $-0.00234$ \\
Random / hybrid hard undersampling 3:1$\dagger$ & 0.43939 / 0.39353 & $-0.00141$ / $-0.04727$ \\
\bottomrule
\end{tabularx}
\end{center}
\vspace{-4pt}

\textbf{Missing values and representation.} Selective imputation combines
Naive Bayes for three categorical columns, ridge for height/weight and BMI
reconstruction. Screening F1 0.44267 uses a threshold selected on the same
predictions being evaluated; inner threshold selection gives 0.44179.
We retain direct missing-value handling. One-hot, daily frequencies and
exact bins also score below the control. Selecting 100 or 60 features
worsens every fold, supporting all 295 columns.

\textbf{Complexity and positive recovery.} More trees, greater depth and
positive weighting fail to improve the historical control. Weights 1.25/1.5
lose F1 in every fold; weight 2 scores 0.43912, adding 126 true positives and
850 false positives.
Two-seed averaging also loses in every fold. The ReLU MLPs (64/32 and 128/64)
use sigmoid output, Adam (0.001, batch 512), binary cross-entropy and L2
$10^{-4}$, with imputation, scaling, one-hot and frequency conversion.
Epoch selection uses an inner 10\% split (maximum 80, patience 8), then refitting.
Both lose in every fold (F1 0.42445/0.41927).

\textbf{Training balance.} Undersampling keeps all positives and 3, 5 or 7
negatives per positive; statistics and bins use retained rows.
Hybrid 3:1 takes half its negatives randomly and half from a training-only
pilot's highest scores (24 trees, depth 3). Its AP falls from 0.42878 to
0.38423, consistent with reduced negative coverage. These matched comparisons
support retaining all rows and weight 1.

\heading{Interpreting the result}
The Phase 34 search adds 206 true positives and 907 false positives
($\Delta$F1 $-0.00135$); AP and log loss barely change. Its paired 95\% interval
[$-0.00300$, $+0.00032$] includes zero. Phase 14 remains the submitted model.
Intervals condition on fixed fits, samples and thresholds.

The selected model finds \textbf{12,855 of 23,180 positives}, with 22,112 false
positives and 10,325 false negatives. At precision 36.76\%, roughly
two thirds of flagged cases are false alarms.

\textbf{HAREHAB1 contributes to the score.} It asks about rehabilitation
after a heart attack; all 620 recorded development responses belong to
positives. Its historical Phase 14 ablation lowers F1 by 0.01112. We retained it to classify
reported disease in this dataset. For prediction before an event, we would
exclude it and check when the other responses would become available.

\heading{Reproduction and limitations}
Results are recorded in \path{results/experiments/}. After
\code{python -m tools.prepare\_data}, the selected CV is reproduced with:\par
{\small\ttfamily python -m src.run\_boosting\_cv configs/experiments/final\_phase14\_validation.json}\par
For final calibration, three stratified OOF folds cover all 328,135 labelled
rows, fitting preprocessing and models within each training fold.
Pooled predictions select threshold \textbf{0.2106169795}. Calibration F1 is
0.442770 for Phase 14 and 0.443570 for Phase 34; these scores use the labels
that selected their thresholds and are not independent evaluations.
\path{configs/final_model.json} fixes the selected implementation and threshold.
The following command trains on all labelled rows:\par
{\small\ttfamily python run.py -{}-data-dir data/raw/dataset -{}-output submission.csv}\par
The verified CSV contains 109,379 predictions: 14,746 positive and 94,633
negative. AIcrowd returned \textbf{F1 0.452, accuracy 0.878} for Phase 14.

\textbf{Development and leaderboard results follow model selection.}
We reused these folds and consulted the local 20\% in initial baselines.
Inner threshold selection reduces one source of optimism. We compared the
two predefined split-search versions on AIcrowd and retained Phase 14;
leaderboard feedback also contributed to the final choice.

\end{document}
